\chapter{Introduction}

\section{Why this matters}

Schools in England continue to face a shortage of teachers in primary and in most secondary
subjects \citep{mclean2025teacher}. As a result, schools rely more heavily on cover and supply
teachers and on late hires, both of which are likely to harm pupil learning
\citep{benhenda2022absence,papay2016productivity}.

Existing research on why teachers leave the profession has focused overwhelmingly on
early-career attrition. This matters for two reasons. First, more teachers leave the profession
during their 30s than during their 20s --- 20\% more, by our count --- so reducing attrition in
this group could materially reduce shortages. Second, teachers leaving in their 30s appear to be
leaving for different reasons than those leaving early in their careers. Those who leave in this
age bracket often cite childcare responsibilities \citep{kersaint2007why} and a lack of flexible
working opportunities \citep{dfe2024working}.

Despite survey evidence suggesting that flexible working would improve retention
\citep{dfe2024working}, a recent review found a lack of experimental or quasi-experimental
evidence on the effects of flexible working in schools \citep{harland2023review}. This project
addresses that gap directly: we use a job choice survey experiment to quantify, for the first
time, how much a wide range of family-friendly working arrangements actually matter to teachers'
decisions to stay in or leave the profession.

\section{Theoretical framing}

We draw on \citet{greenhaus1985sources}'s account of work-family conflict, which distinguishes
two mechanisms through which work and family roles become incompatible: \emph{time-based
conflict}, where time devoted to one role limits participation in another (for example, long
working hours leaving too little time for childcare), and \emph{strain-based conflict}, where
strain generated in one role spills over into the other (for example, workplace stress affecting
home life). The family-friendly working arrangements we test are chosen to address one or both
of these mechanisms.

\section{What we already know, and what we don't}

A small number of job choice survey experiments have now been conducted with in-service teachers
\citep{johnston2021preferences,lentini2024teachers,levatino2024unveiling,lovison2024investing},
including two using English data \citep{allen2025evaluating,burge2021understanding}. These
studies establish that job choice experiments can credibly recover teachers' preferences, but
each looks at a limited range of family-friendly arrangements, and none reports the kind of
detailed subgroup analysis --- by age, gender, and parenting responsibilities --- needed to
understand \emph{why} teachers in their 30s in particular are leaving. \citet{lovison2024investing},
for example, look at workload and childcare subsidies but only report subgroups based on very
wide age brackets; \citet{burge2021understanding} look at workload and part-time working but only
split their sample above/below age 35; \citet{allen2025evaluating} test highly subsidised
childcare and general ``flexible working,'' comparing parents against non-parents, but that study
is designed to understand how schools attract teachers \emph{already in the profession} from
other schools, not to understand what would stop teachers leaving the profession altogether.

This project goes beyond the existing literature in three ways. First, we test nine
family-friendly working attributes not covered in any single existing study, giving by far the
most extensive assessment to date of what matters for retaining experienced teachers. Second,
we frame our job attributes in sector-neutral language that applies equally to teaching and
non-teaching jobs, allowing us to model the decision to leave the profession rather than only the
decision to move between schools. Third, we go beyond comparing the value teachers place on each
arrangement against its cost to schools \emph{on average}, and instead compare it specifically for
the subgroup of teachers who are parents --- the group most likely to be affected by work-family
conflict, and the group whose retention this project is ultimately about.

\section{Research questions}

\begin{description}
  \item[RQ1] What is the relative importance of family-friendly working arrangements, compared
    to salary and other workplace characteristics, in determining teachers' job choices?
  \item[RQ2] What salary-equivalent value do teachers place on these family-friendly working
    arrangements, and how does this compare to the cost of providing them?
  \item[RQ3] How do teachers' preferences for different job attributes vary by (a) age, (b)
    gender, (c) parenting responsibilities, and (d) stated intention to quit teaching?
\end{description}

% NOTE (quick-and-dirty draft, [date]): citation keys above (mclean2025teacher,
% benhenda2022absence, etc.) are placeholders matching the bibliography style used
% in the bid — add matching entries to references.bib before this compiles cleanly
% with citations rather than just placeholder text. Full reference list is in the
% bid document ("Final submission bid as submitted") in the Nuffield bid project.
